% GENERATED FILE. Edit canonical lessons, then run npm run generate:books.
% canonical-source-hash: fnv1a64:a8d9040295a9f2e0
% canonical-lessons: KA-C376-purave, KA-C376-sabitupadisu, KA-C376-sadhyate, KA-C376-mahatva, KA-C376-adharisu

\chapter{Proof and Possibility}
\label{ch:ka-proof-and-possibility}
\hlchaptermodality{376}

\begin{chapteropening}
\textbf{By the end of this chapter:} \emph{I can say proof, to prove, a possibility, importance and to base on.}

Complete the last lesson of chapter 376: I can say proof, to prove, a possibility, importance and to base on.
\end{chapteropening}

\section[\texorpdfstring{\kn{ಪುರಾವೆ} (purāve)}{ಪುರಾವೆ (purāve)}]{\kn{ಪುರಾವೆ} (purāve) — proof; evidence}

\label{lesson:KA-C376-purave}



\begin{quote}
Before the new one: say the Kannada for chalk, then the Kannada for a ruler.
\end{quote}

\subsection*{You'll want to know: \kn{ಪುರಾವೆ}}
\textbf{\kn{ಪುರಾವೆ}} — \emph{purāve} — \textquotedblleft{}proof; evidence\textquotedblright{}.

\begin{quote}
\kn{ಅವನ} \kn{ಮಾತಿಗೆ} \kn{ಯಾವುದೇ} \kn{ಪುರಾವೆ} \kn{ಇರಲಿಲ್ಲ}.
\end{quote}

\begin{quote}
\emph{avana mātige yāvudē purāve iralilla.}
\end{quote}

\begin{quote}
There was no proof for what he said.
\end{quote}

With \kn{ಇದೆ} and \kn{ಇಲ್ಲ}: \kn{ಪುರಾವೆ} \kn{ಇದೆ}, there is proof.

\begin{cousinweb}
\textbf{\kn{ಪುರಾವೆ}} (\emph{purāve}) has no entry in DEDR or CDIAL, so its origin is not settled. Marathi uses a close form, \emph{purāvā}, for evidence.
\end{cousinweb}

\begin{grammarlens}[title={What you've built}]
One more word for giving reasons and opinions.
\end{grammarlens}

\subsection*{Your turn}
\begin{itemize}
\raggedright
  \item \emph{Say it:} \emph{purāve}
  \item \emph{Say it:} \emph{purāve}, once more
  \item \emph{Say it:} \emph{avana mātige yāvudē purāve iralilla.} — There was no proof for what he said.
  \item \emph{Recall:} say the Kannada for chalk, then the Kannada for a ruler, then say \emph{purāve} again
\end{itemize}

\begin{tcolorbox}[breakable,colback=teal!4,colframe=teal!35!black,title={Before you move on}]
What does \kn{ಪುರಾವೆ} mean? (\textquotedblleft{}Proof; evidence\textquotedblright{}.) Say it once more.
\end{tcolorbox}
\section[\texorpdfstring{\kn{ಸಾಬೀತುಪಡಿಸು} (sābītupaḍisu)}{ಸಾಬೀತುಪಡಿಸು (sābītupaḍisu)}]{\kn{ಸಾಬೀತುಪಡಿಸು} (sābītupaḍisu) — to prove}

\label{lesson:KA-C376-sabitupadisu}



\begin{quote}
Before the new one: say the Kannada for a ruler, then the Kannada for proof.
\end{quote}

\subsection*{You'll want to know: \kn{ಸಾಬೀತುಪಡಿಸು}}
\textbf{\kn{ಸಾಬೀತುಪಡಿಸು}} — \emph{sābītupaḍisu} — \textquotedblleft{}to prove\textquotedblright{}.

\begin{quote}
\kn{ನನ್ನ} \kn{ಮಾತು} \kn{ನಿಜ} \kn{ಅಂತ} \kn{ನಾನು} \kn{ಸಾಬೀತುಪಡಿಸಿದೆ}.
\end{quote}

\begin{quote}
\emph{nanna mātu nija anta nānu sābītupaḍiside.}
\end{quote}

\begin{quote}
I proved that my words were true.
\end{quote}

Past (I …): \kn{ಸಾಬೀತುಪಡಿಸಿದೆ}, \emph{sābītupaḍiside}.

\kn{ಸಾಬೀತು} means proved; -\kn{ಪಡಿಸು} makes a verb of it. \kn{ಸಾಬೀತಾಯಿತು}, it was proved.

\begin{cousinweb}
\textbf{\kn{ಸಾಬೀತುಪಡಿಸು}} (\emph{sābītupaḍisu}) is a loan from Hindustani \emph{sābit}, \textquotedblleft{}proved, established\textquotedblright{}, which comes from Arabic \emph{thābit}, \textquotedblleft{}firm\textquotedblright{}, with the Kannada ending -\emph{paḍisu} that makes a verb.
\end{cousinweb}

\begin{grammarlens}[title={What you've built}]
One more word for giving reasons and opinions.
\end{grammarlens}

\subsection*{Your turn}
\begin{itemize}
\raggedright
  \item \emph{Say it:} \emph{sābītupaḍisu}
  \item \emph{Say it:} \emph{sābītupaḍisu}, once more
  \item \emph{Say it:} \emph{nanna mātu nija anta nānu sābītupaḍiside.} — I proved that my words were true.
  \item \emph{Recall:} say the Kannada for a ruler, then the Kannada for proof, then say \emph{sābītupaḍisu} again
\end{itemize}

\begin{tcolorbox}[breakable,colback=teal!4,colframe=teal!35!black,title={Before you move on}]
What does \kn{ಸಾಬೀತುಪಡಿಸು} mean? (\textquotedblleft{}To prove\textquotedblright{}.) Say it once more.
\end{tcolorbox}
\section[\texorpdfstring{\kn{ಸಾಧ್ಯತೆ} (sādhyate)}{ಸಾಧ್ಯತೆ (sādhyate)}]{\kn{ಸಾಧ್ಯತೆ} (sādhyate) — a possibility; a chance}

\label{lesson:KA-C376-sadhyate}



\begin{quote}
Before the new one: say the Kannada for proof, then the Kannada for to prove.
\end{quote}

\subsection*{You'll want to know: \kn{ಸಾಧ್ಯತೆ}}
\textbf{\kn{ಸಾಧ್ಯತೆ}} — \emph{sādhyate} — \textquotedblleft{}a possibility; a chance\textquotedblright{}.

\begin{quote}
\kn{ನಾಳೆ} \kn{ಮಳೆ} \kn{ಬರುವ} \kn{ಸಾಧ್ಯತೆ} \kn{ಇದೆ}.
\end{quote}

\begin{quote}
\emph{nāḷe maḷe baruva sādhyate ide.}
\end{quote}

\begin{quote}
There is a chance of rain tomorrow.
\end{quote}

A verb in -\kn{ುವ} comes before it: \kn{ಬರುವ} \kn{ಸಾಧ್ಯತೆ}, the chance of coming.

\begin{cousinweb}
\textbf{\kn{ಸಾಧ್ಯತೆ}} (\emph{sādhyate}) is borrowed from Sanskrit \emph{sādhyatā}, \textquotedblleft{}feasibility\textquotedblright{} (Monier-Williams), built on \emph{sādhya}, \textquotedblleft{}that can be done\textquotedblright{}, with the ending -\emph{tā}, \textquotedblleft{}-ness\textquotedblright{}.
\end{cousinweb}

\begin{grammarlens}[title={What you've built}]
One more word for giving reasons and opinions.
\end{grammarlens}

\subsection*{Your turn}
\begin{itemize}
\raggedright
  \item \emph{Say it:} \emph{sādhyate}
  \item \emph{Say it:} \emph{sādhyate}, once more
  \item \emph{Say it:} \emph{nāḷe maḷe baruva sādhyate ide.} — There is a chance of rain tomorrow.
  \item \emph{Recall:} say the Kannada for proof, then the Kannada for to prove, then say \emph{sādhyate} again
\end{itemize}

\begin{tcolorbox}[breakable,colback=teal!4,colframe=teal!35!black,title={Before you move on}]
What does \kn{ಸಾಧ್ಯತೆ} mean? (\textquotedblleft{}A possibility; a chance\textquotedblright{}.) Say it once more.
\end{tcolorbox}
\section[\texorpdfstring{\kn{ಮಹತ್ವ} (mahatva)}{ಮಹತ್ವ (mahatva)}]{\kn{ಮಹತ್ವ} (mahatva) — importance}

\label{lesson:KA-C376-mahatva}



\begin{quote}
Before the new one: say the Kannada for to prove, then the Kannada for a possibility.
\end{quote}

\subsection*{You'll want to know: \kn{ಮಹತ್ವ}}
\textbf{\kn{ಮಹತ್ವ}} — \emph{mahatva} — \textquotedblleft{}importance\textquotedblright{}.

\begin{quote}
\kn{ಈ} \kn{ಸಭೆಗೆ} \kn{ತುಂಬಾ} \kn{ಮಹತ್ವ} \kn{ಇದೆ}.
\end{quote}

\begin{quote}
\emph{ī sabhege tumbā mahatva ide.}
\end{quote}

\begin{quote}
This meeting is very important.
\end{quote}

Before a noun it takes -\kn{ದ}: \kn{ಮಹತ್ವದ} \kn{ವಿಷಯ}, an important matter.

\begin{cousinweb}
\textbf{\kn{ಮಹತ್ವ}} (\emph{mahatva}) is borrowed from Sanskrit \emph{mahattva}, \textquotedblleft{}greatness\textquotedblright{} (CDIAL 9933), from \emph{mahat}, \textquotedblleft{}great\textquotedblright{}, with -\emph{tva}, \textquotedblleft{}-ness\textquotedblright{}.
\end{cousinweb}

\begin{grammarlens}[title={What you've built}]
One more word for giving reasons and opinions.
\end{grammarlens}

\subsection*{Your turn}
\begin{itemize}
\raggedright
  \item \emph{Say it:} \emph{mahatva}
  \item \emph{Say it:} \emph{mahatva}, once more
  \item \emph{Say it:} \emph{ī sabhege tumbā mahatva ide.} — This meeting is very important.
  \item \emph{Recall:} say the Kannada for to prove, then the Kannada for a possibility, then say \emph{mahatva} again
\end{itemize}

\begin{tcolorbox}[breakable,colback=teal!4,colframe=teal!35!black,title={Before you move on}]
What does \kn{ಮಹತ್ವ} mean? (\textquotedblleft{}Importance\textquotedblright{}.) Say it once more.
\end{tcolorbox}
\section[\texorpdfstring{\kn{ಆಧರಿಸು} (ādharisu)}{ಆಧರಿಸು (ādharisu)}]{\kn{ಆಧರಿಸು} (ādharisu) — to base on}

\label{lesson:KA-C376-adharisu}



\begin{quote}
Before the new one: say the Kannada for a possibility, then the Kannada for importance.
\end{quote}

\subsection*{You'll want to know: \kn{ಆಧರಿಸು}}
\textbf{\kn{ಆಧರಿಸು}} — \emph{ādharisu} — \textquotedblleft{}to base on\textquotedblright{}.

\begin{quote}
\kn{ಈ} \kn{ಸಿನಿಮಾ} \kn{ಒಂದು} \kn{ನಿಜ} \kn{ಘಟನೆಯನ್ನು} \kn{ಆಧರಿಸಿದೆ}.
\end{quote}

\begin{quote}
\emph{ī sinimā ondu nija ghaṭaneyannu ādhariside.}
\end{quote}

\begin{quote}
This film is based on a true incident.
\end{quote}

Past (I …): \kn{ಆಧರಿಸಿದೆ}, \emph{ādhariside}.

What it rests on takes -\kn{ಅನ್ನು}: \kn{ಘಟನೆಯನ್ನು} \kn{ಆಧರಿಸಿ}, based on an incident.

\begin{cousinweb}
\textbf{\kn{ಆಧರಿಸು}} (\emph{ādharisu}) is built on Sanskrit \emph{ādhāra}, \textquotedblleft{}a support, a basis\textquotedblright{} (CDIAL 1165), with the Kannada ending -\emph{isu} that turns a Sanskrit noun into a verb.
\end{cousinweb}

\begin{grammarlens}[title={What you've built}]
One more word for giving reasons and opinions.
\end{grammarlens}

\subsection*{Your turn}
\begin{itemize}
\raggedright
  \item \emph{Say it:} \emph{ādharisu}
  \item \emph{Say it:} \emph{ādharisu}, once more
  \item \emph{Say it:} \emph{ī sinimā ondu nija ghaṭaneyannu ādhariside.} — This film is based on a true incident.
  \item \emph{Recall:} say the Kannada for a possibility, then the Kannada for importance, then say \emph{ādharisu} again
\end{itemize}

\begin{tcolorbox}[breakable,colback=teal!4,colframe=teal!35!black,title={Before you move on}]
What does \kn{ಆಧರಿಸು} mean? (\textquotedblleft{}To base on\textquotedblright{}.) Say it once more.
\end{tcolorbox}
